% SPDX-FileCopyrightText: 2026 Will Cook
% SPDX-License-Identifier: CC-BY-4.0
\documentclass[11pt]{article}
\usepackage[T1]{fontenc}
\usepackage[a4paper,margin=1in]{geometry}
\usepackage[dvipsnames]{xcolor}
\usepackage[colorlinks=true,urlcolor=RoyalBlue]{hyperref}
\hypersetup{
  pdftitle={SCRAPPED — Former Erdős Problem 249 specialist note},
  pdfauthor={Will Cook},
  pdfsubject={Retirement notice for the former one-sided specialist note},
}
\begin{document}
\begin{center}
{\Huge\bfseries\color{BrickRed} SCRAPPED}\par
\vspace{1em}
{\Large Former Erd\H{o}s Problem 249 specialist note}\par
\vspace{0.8em}
{\large This is not an active mathematical paper.}
\end{center}

\section*{Why this note was retired}

The repository studies two Erd\H{o}s problems, 249 and 257. Publishing a
standalone technical companion for only Problem 249 gives one exploratory
route a status that the overall paper programme does not justify. It also
fragments the reader's route through mathematics that is already being
reworked in the main exposition.

The former note is therefore scrapped as a publication artefact. Its
mathematics has not been declared false; the problem is the one-sided paper
boundary and the resulting duplication.

\section*{The only two coherent successors}

\begin{enumerate}
  \item Keep one main mathematical paper and fold only the necessary
  specialist exposition into it; or
  \item revive a matched pair of specialist notes, one for Problem 249 and
  one for Problem 257, with equally clear jobs and boundaries.
\end{enumerate}

There is no active plan to polish the former 249-only note in isolation.

\section*{What was preserved}

The full source and rendered PDF that preceded this notice remain in
\texttt{paper/archive/} under filenames ending
\texttt{-retired-2026-07-17}. They are retained as provenance and a source of
potentially reusable exposition, not as current publications.

The main Erd\H{o}s mathematics paper remains a separate live manuscript. This
retirement notice makes no mathematical status change and must not be used to
infer that either open problem has been solved.

\vfill
\noindent Retired 17 July 2026.
\end{document}
